\section{Dashboard de observabilidad}

El sistema expone un panel web de una sola pagina que muestra, en tiempo real,
el estado de cada nodo del cluster: rol, numero de cuentas, saldo total, numero
de transferencias, identificador de la ultima transaccion, porcentajes de CPU,
RAM y disco, y el numero de transacciones registradas en Cloud Storage. El panel
no necesita servidor adicional ni dependencias de frontend: es HTML estatico
servido por el propio nodo y alimentado por dos endpoints JSON.

\subsection{Pagina servida en GET /}

El recurso \codigo{src/main/resources/dashboard.html} se empaqueta dentro del jar
y lo entrega \codigo{PaginaHandler} en \codigo{GET /}. Para evitar leer el archivo
en cada peticion, el manejador lo carga una sola vez en un arreglo de bytes
estatico y lo cachea; si el recurso falta responde 500 y a cualquier ruta distinta
de la raiz responde 404.

\begin{lstlisting}[language=Java,caption={PaginaHandler.java: dashboard cacheado en bytes}]
private static final byte[] HTML = cargar();

@Override
public Respuesta manejar(Solicitud s) {
    if (!"GET".equalsIgnoreCase(s.metodo())) return Respuesta.sinCuerpo(405);
    if (!"/".equals(s.path())) return Respuesta.sinCuerpo(404);
    if (HTML == null) return Respuesta.sinCuerpo(500);
    return Respuesta.html(200, HTML);
}
\end{lstlisting}

\subsection{Estado de un nodo: GET /stats}

\codigo{StatsHandler} responde en \codigo{GET /stats} con el snapshot del nodo
local. Es publico (sin JWT) porque solo expone monitoreo de un nodo. El metodo
estatico \codigo{StatsHandler.snapshot()} construye un mapa ordenado que tambien
reutiliza el panel agregado. Determina el rol leyendo la variable de entorno
\codigo{BANCO\_LIDER\_HOST} (vacia implica \codigo{LIDER}, definida implica
\codigo{REPLICA}) y publica el saldo total en centavos como \codigo{long} exacto
para verificar el invariante de consistencia en una sola peticion.

\begin{lstlisting}[language=Java,caption={StatsHandler.java: snapshot del nodo}]
out.put("rol", rol);
out.put("cuentas", repo.tamano());
out.put("saldoTotalCentavos", totalCent);
out.put("transferencias", repo.totalTransferencias());
out.put("ultimaTxId", repo.ultimaTxId());
out.put("secuencia", repo.secuenciaActual());
out.put("cpu", SistemaMetricas.cpuPorcentaje());
out.put("ram", SistemaMetricas.ramPorcentaje());
out.put("disco", SistemaMetricas.discoPorcentaje());
out.put("txEnStorage", repo.txEnStorage()); // -1 si el log GCS esta deshabilitado
\end{lstlisting}

Las metricas de sistema viven en \codigo{SistemaMetricas} y se calculan solo con
APIs del JDK, sin dependencias externas: el \codigo{\%}CPU y el \codigo{\%}RAM se
obtienen del \codigo{OperatingSystemMXBean} de \codigo{com.sun.management}, y el
\codigo{\%}Disco se mide con un \codigo{File} sobre la raiz del sistema de
archivos.

\begin{lstlisting}[language=Java,caption={SistemaMetricas.java: CPU y RAM sin dependencias}]
public static double cpuPorcentaje() {
    double carga = OS.getCpuLoad();
    return carga < 0 ? 0.0 : redondear(carga * 100.0);
}

public static double ramPorcentaje() {
    long total = OS.getTotalMemorySize();
    if (total <= 0) return 0.0;
    long usada = total - OS.getFreeMemorySize();
    return redondear(100.0 * usada / total);
}
\end{lstlisting}

\subsection{Vista agregada del cluster: GET /panel}

\codigo{PanelHandler} responde en \codigo{GET /panel} y construye la vista
completa del cluster que consume el dashboard. El lider toma su propio
\codigo{StatsHandler.snapshot()} y, ademas, consulta el \codigo{/stats} de cada
peer. La consulta es un \codigo{fetch} del lado del servidor (con
\codigo{HttpClient}), lo que evita problemas de CORS en el navegador y centraliza
la agregacion en un solo origen. Cada nodo se envuelve con metadatos
\codigo{host} y \codigo{alcanzable}; si un peer no responde dentro del timeout de
2 s, se marca \codigo{alcanzable:false} sin abortar el resto del panel.

\begin{lstlisting}[language=Java,caption={PanelHandler.java: agregado lider + peers}]
List<Map<String, Object>> nodos = new ArrayList<>();
nodos.add(conMeta(idLocal, true, StatsHandler.snapshot()));
for (String peer : peers) nodos.add(consultarPeer(peer));
// ...
private Map<String, Object> consultarPeer(String hostPuerto) {
    try {
        HttpRequest req = HttpRequest.newBuilder(URI.create("http://" + hostPuerto + "/stats"))
                .timeout(Duration.ofSeconds(2)).GET().build();
        String cuerpo = http.send(req, BodyHandlers.ofString()).body();
        // ...
        return conMeta(hostPuerto, true, m);
    } catch (Exception e) {
        return conMeta(hostPuerto, false, new LinkedHashMap<>());
    }
}
\end{lstlisting}

\begin{nota}
Como \codigo{HttpClient.send} hacia los peers es bloqueante (hasta 2 s),
\codigo{PanelHandler} se asigna al pool WORKER; si corriera en un reactor lo
estancaria. Asi un peer caido nunca cuelga el panel del resto del cluster.
\end{nota}

\subsection{Refresco reactivo en el navegador}

El JavaScript embebido en \codigo{dashboard.html} consulta \codigo{/panel},
dibuja una tarjeta por nodo (con badge \codigo{LIDER}, \codigo{REPLICA} o
\codigo{CAIDO} segun el campo \codigo{alcanzable}) y compara los
\codigo{saldoTotalCentavos} de los nodos vivos para senalar si los saldos son
consistentes entre nodos. El refresco es reactivo: tras la carga inicial, un
\codigo{setInterval} repite la consulta cada 2 segundos, cubriendo el punto
extra de actualizacion automatica.

\begin{lstlisting}[language=none,caption={dashboard.html: refresco automatico cada 2 s}]
function refrescar() {
  fetch("/panel").then(function(r){ return r.json(); })
    .then(pintar)
    .catch(function(){ document.getElementById("total").innerHTML =
      '<span class="mal">no se pudo contactar al nodo</span>'; })
    .finally(function(){
      document.getElementById("reloj").textContent = new Date().toLocaleTimeString("es-MX");
    });
}

refrescar();
setInterval(refrescar, 2000);
\end{lstlisting}
